\section{Limitations}

\subsection{Decoder Fallback}

For decoders with $L > k$, we fall back to dense attention. This **negates 
the linear complexity benefit** for the case where it's most needed.

\subsection{Parameter Increase}

Linformer adds ~7\% parameters per layer. This is a real cost.

\subsection{Sequence Length Extrapolation}

For $L > L_\text{max}$, we interpolate $E, F$. This is a heuristic not 
validated on very long sequences.

\subsection{Task Coverage}

Only SST-2 evaluated. Generation, multi-class, and NER tasks untested.

\subsection{Model Coverage}

Tested up to 124M params. Behavior on larger models unknown.

\subsection{Single-Seed Results}

SE ~1.7\%. Differences below 3.4\% may be noise.

\subsection{No Comparison with Original Linformer}

We do not reproduce Wang et al. (2020) exactly. Our diagonal init, adaptive 
slicing, and causal fallback differ.

\subsection{E, F Memory}

$E, F \in \mathbb{R}^{k \times L_\text{max}}$ cost $2 k L_\text{max}$ 
parameters per layer. For large $L_\text{max}$, this grows.

\subsection{Causal Analysis}

Our causal fallback is a \textbf{correctness} measure, not a \textbf{performance} 
measure. It sacrifices efficiency.

\subsection{The output\_scale Parameter}

\texttt{output\_scale} is a scalar initialized to 1. It's trainable but rarely 
moves significantly.

\subsection{Code Quality}

Research-grade:
\begin{itemize}
\item Broad exception handling
\item No fused kernel
\item Some numerical stability concerns

\end{itemize}
\subsection{Ethical Considerations}

Dual-use. See Appendix E.
\begin{lstlisting}

---
\end{lstlisting}